\section{Evaluation}

Evaluating our approach requires a reference dataset that establishes an explicit link between a textual description of a business rule and the piece of source code implementing that rule. Such a link is essential to determine whether the proposed approach is able to retrieve the correct implementation among the candidate methods returned for a given business rule. More specifically, for each business rule used during the evaluation, we need to know the method that actually implements it, so that the presence or absence of this method in the final set of retrieved candidates can be assessed.

However, establishing such links is challenging in the context of legacy software systems. In practice, source code and business documentation are often maintained as separate artifacts. Although documentation may contain detailed descriptions of business rules, it generally does not explicitly indicate where these rules are implemented in the source code. Conversely, the source code may contain the implementation of a business rule without providing a direct reference to its corresponding documentation. Consequently, the lack of traceability between business rules and their implementations makes it difficult to construct a reliable ground-truth dataset for evaluating rule-location approaches.

This challenge is particularly relevant in our case study. We have access to the source code of the studied application and to a collection of documents containing descriptions of business rules. However, these documents do not provide explicit links to the methods implementing the corresponding rules. Therefore, we had to manually reconstruct these links in order to build an evaluation dataset in the form of business rule–implementation pairs.

To guide this manual identification process, we relied primarily on the information available within the source code. In particular, source-code comments and method names provided valuable clues for identifying potential implementations. Some comments contain explicit references to business-rule identifiers corresponding to the numbering used in the documentation. These references allowed us to establish a connection between a documented business rule and a specific part of the source code. We systematically inspected such occurrences and manually verified whether the referenced code actually implemented the corresponding business rule.

Through this manual investigation, we were able to identify 100 business rules for which a corresponding implementation method could be established with sufficient confidence. These pairs constitute our reference dataset, with each instance associating a business rule description from the documentation with the source-code method implementing that rule. This dataset is subsequently used as the ground truth for evaluating the ability of our approach to retrieve the correct implementation among the candidate methods.